\chapter{準備}
\label{ch:preliminaries}

本章では，最適輸送の定式化に必要な概念を導入する．

\section{$\sigma$-代数}
\label{sec:sigma-algebra}

\begin{definition}{$\sigma$-代数}{amb-sigma-algebra}
  集合 $X$ の部分集合族 $\mathcal{F} \subset 2^X$ が
  \textbf{$\sigma$-代数}であるとは，次の3条件を満たすことをいう：
  \begin{enumerate}
    \item $X \in \mathcal{F}$．
    \item $A \in \mathcal{F} \implies X \setminus A \in \mathcal{F}$（補集合で閉じる）．
    \item $A_1, A_2, \ldots \in \mathcal{F} \implies
      \bigcup_{i=1}^\infty A_i \in \mathcal{F}$（可算合併で閉じる）．
  \end{enumerate}
\end{definition}

\section{Borel 集合と Borel 測度}
\label{sec:borel}

\begin{definition}{Borel $\sigma$-代数}{amb-borel-sigma}
  $X \subset \R^n$ とする．$X$ の開集合全体を含む最小の $\sigma$-代数を
  $X$ の\textbf{Borel $\sigma$-代数}といい，$\Borel(X)$ と書く．
  $\Borel(X)$ の元を\textbf{Borel 集合}という．
\end{definition}

\begin{definition}{Borel 測度}{amb-borel-measure}
  $X \subset \R^n$ 上の\textbf{Borel 測度}とは，
  $\Borel(X)$ 上で定義された測度
  $\mu \colon \Borel(X) \to [0, \infty]$ のことをいう．
\end{definition}

\section{Lebesgue 測度}
\label{sec:lebesgue}

\begin{definition}{Lebesgue 測度}{amb-lebesgue}
  $\R^n$ 上の\textbf{Lebesgue 測度} $\Leb^n$ とは，
  区間の直積 $\prod_{i=1}^n [a_i, b_i]$ に対し
  \[
    \Leb^n\!\left(\prod_{i=1}^n [a_i, b_i]\right) = \prod_{i=1}^n (b_i - a_i)
  \]
  を満たす $\Borel(\R^n)$ 上の測度である．
\end{definition}

Lebesgue 測度は Borel 測度だが，$\Leb^n(\R^n) = \infty$ であり，
有限測度ではない．

\section{Radon 測度}
\label{sec:radon}

\begin{definition}{Radon 測度}{amb-radon}
  $X \subset \R^n$ 上の Borel 測度 $\mu$ が\textbf{Radon 測度}であるとは，
  次の2条件を満たすことをいう：
  \begin{enumerate}
    \item \textbf{局所有限性}：$\forall$ コンパクト集合 $K \subset X$ に対し $\mu(K) < \infty$．
    \item \textbf{内部正則性}：$\forall\, B \in \Borel(X)$ に対し
      \[
        \mu(B) = \sup\,\{\,\mu(K) \mid K \subset B,\; K \text{ はコンパクト}\,\}.
      \]
  \end{enumerate}
\end{definition}

\begin{example}{Lebesgue 測度は Radon 測度}{amb-lebesgue-radon}
  $\Leb^n$ は Radon 測度である．
  $\forall$ コンパクト集合 $K \subset \R^n$ は有界だから $\Leb^n(K) < \infty$（局所有限性）．
  内部正則性も成り立つ．
\end{example}

\section{確率測度}
\label{sec:probability}

\begin{definition}{確率測度}{amb-probability}
  $X \subset \R^n$ 上の Borel 測度 $\mu$ が\textbf{確率測度}であるとは，
  $\mu \geq 0$ かつ $\mu(X) = 1$ を満たすことをいう．
  $X$ 上の確率測度の全体を $\Prob(X)$ と書く．
\end{definition}

\section{双対空間}
\label{sec:dual}

\begin{definition}{有界線形汎関数と双対空間}{amb-dual}
  ノルム空間 $V$ 上の線形写像 $\Lambda \colon V \to \R$ が
  \textbf{有界}であるとは，
  \[
    \|\Lambda\| := \sup_{\|v\| \leq 1} |\Lambda(v)| < \infty
  \]
  を満たすことをいう．
  $V$ 上の有界線形汎関数の全体を\textbf{双対空間}といい，$V^*$ と書く．
\end{definition}

\section{Riesz の表現定理}
\label{sec:riesz}

\begin{theorem}{Riesz の表現定理}{amb-riesz}
  $X$ をコンパクト集合とする．
  $\forall\, \Lambda \in C(X)^*$ に対し，$\exists!\, \mu \in \Meas(X)$ で
  \[
    \Lambda(f) = \int_X f\,\d\mu
    \qquad \forall\, f \in C(X)
  \]
  を満たすものが存在する．さらに $\|\Lambda\| = |\mu|(X)$ が成り立つ．
\end{theorem}

\begin{proof}
  $\Lambda \geq 0$（$f \geq 0 \implies \Lambda(f) \geq 0$）の場合を示す．

  開集合 $V \subset X$ に対し
  \[
    \mu(V) := \sup\,\{\,\Lambda(f) \mid f \in C(X),\; 0 \leq f \leq 1,\; \spt(f) \subset V\,\}
  \]
  と定め，$\forall$ Borel 集合 $E$ に対し
  $\mu(E) := \inf\,\{\,\mu(V) \mid V \supset E,\; V \text{ は開}\,\}$
  で外側から拡張する．

  $\mu$ が $\sigma$-加法的な Borel 測度であることは，
  $X$ のコンパクト性と $\Lambda$ の正値性・線形性から検証できる．
  内部正則性も $X$ のコンパクト性から従い，$\mu$ は Radon 測度である．

  $\Lambda(f) = \int_X f\,\d\mu$ は，
  $f$ を $0 \leq f \leq 1$ の場合に上の構成から直接確かめ，
  線形性で一般の $f \in C(X)$ に拡張する．

  一般の $\Lambda \in C(X)^*$ は $\Lambda = \Lambda^+ - \Lambda^-$
  （$\Lambda^+(f) := \sup\,\{\,\Lambda(g) \mid 0 \leq g \leq f\,\}$）
  と正の部分に分解でき，各々に上の構成を適用して
  $\mu = \mu^+ - \mu^-$ を得る．
\end{proof}

\section{Banach–Alaoglu の定理}
\label{sec:banach-alaoglu}

\begin{theorem}{Banach–Alaoglu}{amb-banach-alaoglu}
  ノルム空間 $V$ の双対空間 $V^*$ の閉単位球
  $B_{V^*} := \{\Lambda \in V^* \mid \|\Lambda\| \leq 1\}$
  は，$\forall\, v \in V$ に対し $\Lambda \mapsto \Lambda(v)$ が
  連続となる最弱の位相（弱収束の位相）でコンパクトである．
\end{theorem}

\begin{proof}
  $\forall\, v \in V$ に対し $D_v := [-\|v\|,\, \|v\|] \subset \R$ とおく．
  写像
  \[
    \Phi \colon B_{V^*} \to \prod_{v \in V} D_v, \qquad
    \Phi(\Lambda) := \bigl(\Lambda(v)\bigr)_{v \in V}
  \]
  は単射かつ連続である．
  $\prod_{v \in V} D_v$ は Tychonoff の定理によりコンパクトである．

  $\Phi(B_{V^*})$ が $\prod D_v$ の閉部分集合であることを示す．
  $(\ell_v)_{v \in V} \in \overline{\Phi(B_{V^*})}$ とすると，
  $\ell_v$ は $v$ について線形であり（線形性は各点収束の極限で保たれる），
  $|\ell_v| \leq \|v\|$ を満たすから，$\ell \in B_{V^*}$ である．

  よって $\Phi(B_{V^*})$ はコンパクト空間の閉部分集合だからコンパクトであり，
  $\Phi$ は $B_{V^*}$ からその像への同相写像だから，$B_{V^*}$ もコンパクトである．
\end{proof}

\section{測度の弱収束}
\label{sec:weak-convergence}

\begin{definition}{弱収束}{amb-weak-convergence}
  $X$ をコンパクト集合とする．
  $\Prob(X)$ 上の列 $\{\mu_n\}$ が $\mu$ に\textbf{弱収束}するとは，
  \[
    \int_X f\,\d\mu_n \to \int_X f\,\d\mu
    \qquad \forall\, f \in C(X)
  \]
  が成り立つことをいう．
\end{definition}

Riesz の表現定理（定理~\ref{thm:amb-riesz}）により，
弱収束は $C(X)^*$ の収束と同一視できる．

\begin{theorem}{$\Prob(X)$ のコンパクト性}{amb-prob-compact}
  $X$ がコンパクトならば，$\Prob(X)$ は弱収束に関してコンパクトである．
\end{theorem}

\begin{proof}
  Riesz の表現定理より $\Meas(X) \cong C(X)^*$ であり，
  弱収束の位相は $C(X)^*$ 上の位相に一致する．
  Banach–Alaoglu の定理（定理~\ref{thm:amb-banach-alaoglu}）により
  $\{\mu \in \Meas(X) \mid \|\mu\| \leq 1\}$ はコンパクトである．
  $\Prob(X)$ はこの閉部分集合だから，コンパクトである．
\end{proof}

\section{下半連続性}
\label{sec:lsc}

\begin{definition}{下半連続関数}{amb-lsc}
  位相空間 $X$ 上の関数 $f \colon X \to (-\infty, \infty]$ が
  \textbf{下半連続}（lower semicontinuous）であるとは，
  $\forall\, \alpha \in \R$ に対し
  $\{\, x \in X \mid f(x) \leq \alpha \,\}$ が閉集合であることをいう．
\end{definition}

連続関数は下半連続である．
逆は一般に成り立たないが，下半連続関数は連続関数の上限として表せる．

\begin{theorem}{下半連続関数の最小値}{amb-lsc-minimum}
  $X$ がコンパクト位相空間で，$f \colon X \to (-\infty, \infty]$ が
  下半連続ならば，$f$ は $X$ 上で最小値を達成する．
\end{theorem}

\begin{proof}
  $m := \inf_{x \in X} f(x)$ とおく．
  $m = \infty$ なら $X = \emptyset$ だが $X$ はコンパクトなのでこれはない．
  $\forall\, \alpha > m$ に対し $F_\alpha := \{x \in X \mid f(x) \leq \alpha\}$ は
  下半連続性より閉かつ非空であり，$\{F_\alpha\}$ は有限交叉性をもつ．
  $X$ のコンパクト性より $\bigcap_{\alpha > m} F_\alpha \neq \emptyset$ であり，
  この交叉に属する点で $f$ は値 $m$ を達成する．
\end{proof}
